\documentclass[11pt]{article}
\usepackage[T1]{fontenc}
\usepackage[utf8]{inputenc}
\usepackage{lmodern,amsmath,amssymb,amsthm,mathtools,booktabs,longtable}
\usepackage[margin=1in]{geometry}
\usepackage{microtype}
\usepackage[hidelinks]{hyperref}
\newtheorem{theorem}{Theorem}[section]
\newtheorem{proposition}[theorem]{Proposition}
\newtheorem{lemma}[theorem]{Lemma}
\theoremstyle{definition}\newtheorem{definition}[theorem]{Definition}
\theoremstyle{remark}\newtheorem{remark}[theorem]{Remark}
\newcommand{\C}{\mathbb C}\newcommand{\Z}{\mathbb Z}\newcommand{\R}{\mathbb R}
\DeclareMathOperator{\tr}{tr}\DeclareMathOperator{\Rea}{Re}
\newcommand{\status}{\textbf{Research draft. GRH is not proved here.}}
\date{November 2, 2026}

\newtheorem{corollary}[theorem]{Corollary}
\newcommand{\gshift}[2]{#1\langle #2\rangle}
\date{November 3, 2026}
\title{Projective resolutions, surface orders and zeta determinants\\\large A constructive framework and the arithmetic comparison problem}
\author{Research working draft}
\begin{document}
\maketitle
\begin{abstract}
We develop the homological interface in a proposed surface-order approach to the generalized Riemann hypothesis. Forbidden-arrow chains give explicit minimal resolutions, Ext and Tor; permitted-arrow states give a short proof of the weighted Cartan determinant formula. We distinguish these formal invariants from determinant lines and Frobenius zeta functions, correct periodicity, grading and base-change errors, and state a precise conditional arithmetic program. The algebraic results have complete proofs; selected matrix and module components are checked in Lean4. The missing number-field determinant and positive realization remain open. This is not a proof of GRH.
\end{abstract}
\tableofcontents
\section{Objects, conventions and scope}
The starting sources are Schreiber's surface-algebra and surface-order papers \cite{S1,S2}, the generic-resolution discussion \cite{S3}, and the weighted locally gentle determinant formula of Bessenrodt--Holm \cite{BH}. Our purpose is to supply the exact homological constructions needed to use their useful parts in a Hilbert--P\'olya program. The source manuscripts are preserved; corrections below are statements in this new account. Full reading records, pinned versions and artifact hashes accompany the project. Related literature is used with the hypotheses stated here, rather than incorporated as an unqualified equivalence.

Fix a field $k$, a finite quiver $Q=(Q_0,Q_1,s,t)$, and an ideal $I$ generated by a set of length-two paths. We compose paths in the order traveled: $ab$ means first $a$, then $b$. Write $A=kQ/I$, $e_i$ for its vertex idempotents, and $J=A_{>0}$ for the arrow ideal. A path is \emph{permitted} when none of its adjacent pairs belongs to $I$. These paths, including the empty paths $e_i$, form a $k$-basis of $A$. The algebra may be infinite dimensional.

\begin{definition}
The pair $(Q,I)$ is locally gentle if each vertex has at most two incoming and two outgoing arrows, and every arrow has at most one permitted successor, one forbidden successor, one permitted predecessor and one forbidden predecessor. A gentle algebra is additionally finite dimensional. We use right modules, $P_i=e_iA$ and $S_i=P_i/P_iJ$, with arrows of degree one and internal shift $(\gshift{M}{n})_d=M_{d-n}$.
\end{definition}
Thus $\gshift{P_i}{n}$ is generated in degree $n$; square brackets later denote the cohomological shift in a derived category. Completion $\widehat A=\varprojlim A/J^r$, a hereditary normalization of a completed order, and a central special fiber are different objects. A forbidden chain is a composable sequence all of whose adjacent pairs lie in $I$; it is generally zero as a path in $A$. We retain it as an indexing word. Forgetting this distinction makes the resolution below disappear.

The number-field target is initially primitive Dirichlet $L(s,\chi)$. Finite-order Hecke characters over a specified number field, Dedekind zeta functions, and general Artin representations require separate comparison theorems. General Artin holomorphy is not assumed. None of the algebraic results below identifies its determinant with any of these functions.

\section{Minimal projective resolutions and the signed Euler matrix}
For $n\geq1$, let $\mathcal R_n(i)$ be the forbidden chains of length $n$ starting at $i$; set $\mathcal R_0(i)=\{e_i\}$. For a nonempty chain $p=a_1\cdots a_n$, write $p^-=a_1\cdots a_{n-1}$.
\begin{theorem}[Explicit linear resolution]\label{thm:resolution}
For every quadratic monomial algebra as above, a graded projective resolution of $S_i$ is
\[
 \cdots\longrightarrow P_n(i)\xrightarrow{d_n}P_{n-1}(i)
 \longrightarrow\cdots\longrightarrow P_0(i)\longrightarrow S_i\longrightarrow0,
 \qquad P_n(i)=\bigoplus_{p\in\mathcal R_n(i)} \gshift{e_{t(p)}A}{n}.
\]
On the $p$ summand, $d_n(x)=a_nx$ in the $p^-$ summand. The augmentation is the quotient $e_iA\to S_i$. The maps are right $A$-linear, homogeneous of degree zero, and valued in the graded arrow radical. The resolution is minimal in the graded sense. In particular $A$ with its natural path grading is Koszul.
\end{theorem}
\begin{proof}
Each summand is an idempotent projective. The endpoints make left multiplication $a_nx$ a map into the indicated right module; associativity proves right linearity. The source generator has degree $n$, and its image has degree $1+(n-1)=n$. Two successive maps multiply by $a_{n-1}a_n=0$.

The augmentation kernel is the span of all nonempty permitted paths from $i$, the direct sum of $aA$ over arrows $a$ from $i$. This is $\operatorname{im}d_1$. More generally, the kernel of $e_{t(a)}A\xrightarrow{x\mapsto ax}e_{s(a)}A$ is the span of permitted paths whose first arrow $b$ satisfies $ab\in I$. Indeed $ae_{t(a)}=a\ne0$; every other basis path either vanishes after multiplication because of that first pair, or maps to a distinct surviving basis path. The kernel is therefore $\bigoplus_{b:ab\in I}bA$. For a fixed prefix $p^-$, images of distinct last arrows have disjoint nonzero path supports. There can be no cancellation between them. Taking their individual kernels identifies $\ker d_n$ with the image of the summands indexed by $pb$. This proves exactness in every degree without a global-dimension assumption. All entries have positive degree (graded minimality does not identify $J$ with the ordinary Jacobson radical), so applying $\operatorname{Hom}_{\mathrm{gr}\text{-}A}(-,\gshift{S_j}{r})$ kills the differentials. Minimality and linearity follow. The direct sum of the $S_i$ is $A/J$, so its linear resolution is the Koszul property.
\end{proof}

\begin{corollary}[Ext, Tor and inversion]\label{cor:euler}
In ungraded modules, $\dim_k\operatorname{Ext}^n_A(S_i,S_j)$ is the number of forbidden chains of length $n$ from $i$ to $j$. The same number is $\dim_k\operatorname{Tor}^A_n(S_i,{}_jS)$, where ${}_jS=Ae_j/JAe_j$. With the above grading, Ext is supported on the matching internal shift $n$. Define
\[
 C_{ij}(z)=\sum_{d\geq0}\dim_k(e_iA_de_j)z^d,\qquad
 E_{ij}(z)=\sum_{n\geq0}(-1)^n\dim_k\operatorname{Ext}^n_A(S_i,S_j)z^n.
\]
Then $E(z)C(z)=C(z)E(z)=1$ in $M_{Q_0}(\mathbb Z[[z]])$.
\end{corollary}
\begin{proof}
An $A$-map $e_vA\to S_j$ is determined by an element of $S_je_v$, of dimension $\delta_{vj}$. Each differential becomes zero since arrows annihilate the simple. Tensoring with a left simple gives the same calculation. Degree by degree, the augmented resolution has Euler characteristic $S_i$. Its $n$th term contributes $(-1)^nz^nC_{t(p),j}(z)$, so the alternating sum is exactly $(EC)_{ij}=\delta_{ij}$. Each coefficient is finite: $n\leq d$ and $Q$ is finite. Since both matrices have constant term the identity, the formal matrix inverse is unique and two-sided.
\end{proof}

This already signed $E$ is evaluated at $z$, not at $-z$. Other conventions can transpose the matrices, but cannot justify applying the sign twice. The complementary-relation algebra $A^\#=kQ/I^\#$ allows precisely the forbidden pairs of $A$, and $E(z)=C_{A^\#}(-z)$. Its paths encode the quadratic Koszul dual, with an opposite quiver if that convention is used for Yoneda composition. At the chain level Yoneda products are concatenation of forbidden chains (zero when the join is permitted), with the handedness/opposite convention just declared. An explicit comparison-map proof is given in the extended edition; the dimensions and inversion above do not require this product identification. Thus endpoint data and relations, rather than Betti numbers alone, control the product.

For locally gentle $A$, each initial arrow has at most one forbidden continuation. The resolution has at most two arms; they terminate or enter a finite cycle. A repeated ordered arrow pattern yields an eventual tail period after forgetting internal grading. This does not say that $\Omega^pS_i\cong S_i$, or that the displayed period is the least module-isomorphism period.

\section{The weighted Cartan determinant from arrow states}
Give each arrow an independent commuting weight $w_a$ and work over $R=\mathbb Z[[w_a:a\in Q_1]]$. A path has weight the product of its arrows. The weighted Cartan matrix $C_{ij}(\mathbf w)$ sums permitted paths from $i$ to $j$. Finite $Q$ guarantees finitely many terms of each total degree.

Put $n=|Q_0|$, $m=|Q_1|$ and define
\[
 S_{i,a}=\mathbf1_{s(a)=i},\quad T_{a,j}=\mathbf1_{t(a)=j},\quad
 D_{a,b}=\delta_{ab}w_a,
\]
\[
 U_{a,b}=\mathbf1_{t(a)=s(b),\,ab\notin I}w_b,\qquad
 V_{a,b}=\mathbf1_{t(a)=s(b),\,ab\in I}w_b.
\]
Here $S$ is $n\times m$, $T$ is $m\times n$, and $U,V,D$ are $m\times m$. Every composable pair is exactly one kind, so $U+V=TSD$.

\begin{theorem}[Weighted cycle formula, cf.\ Bessenrodt--Holm]\label{thm:cartan}
For a finite locally gentle quiver,
\[
 \det C_A(\mathbf w)=
 \frac{\displaystyle\prod_{c\in\mathcal C_{\rm forb}}
       (1-(-1)^{\ell(c)}w(c))}
      {\displaystyle\prod_{c\in\mathcal C_{\rm perm}}(1-w(c))}.
\]
Both sets consist of primitive directed cycles of the corresponding arrow-transition graph, modulo cyclic rotation. Iterations are not extra factors. All denominators have constant term one.
\end{theorem}
\begin{proof}
Paths of length $r+1$ contribute $SDU^rT$. Hence
\[
 C=1+SD(1-U)^{-1}T,
 \qquad (1-U)^{-1}=\sum_{r\geq0}U^r.
\]
The inverse converges in the augmentation topology because every entry of $U$ has positive degree. Sylvester's determinant identity $\det(1+AB)=\det(1+BA)$ for rectangular $A,B$ gives
\[
 \det C=\det\bigl(1+(1-U)^{-1}TSD\bigr)
        =\frac{\det(1+V)}{\det(1-U)}.
\]
For completeness the rectangular identity follows by expanding the determinant of the block matrix $\left(\begin{smallmatrix}1&-A\\B&1\end{smallmatrix}\right)$ with either identity block. Thus it is valid over any commutative ring. Multiplying before division gives a cleared identity even in more general coefficient rings.

Each vertex of the $U$ transition graph has at most one predecessor and one successor, so every connected component is a directed chain or a directed cycle. The same holds for $V$. Ordering a chain along its arrows makes its transition matrix strictly triangular, contributing determinant one. On a cycle of length $r$, only the identity permutation and the full cyclic permutation contribute to $\det(1-U)$; their sum is $1-w(c)$. Replacing $U$ by $-V$ gives $1-(-1)^rw(c)$. For a one-loop component this is simply $1-w_a$ or $1+w_a$, so there is no exceptional sign case. Multiplying the component determinants proves the formula, including quivers with both infinite permitted paths and infinite resolutions.
\end{proof}

The transfer identity before component factorization holds for every finite quadratic monomial quiver. The locally gentle hypothesis is used for the chain/cycle decomposition, not for Sylvester's identity. Formula~\ref{thm:cartan} is a formal invariant, not the determinant of an infinite resolution. Specializing weights to a point where a denominator vanishes requires a separate cancellation or continuation argument.

\section{Completion, singular seams and geometric objects}
The node illustrates both the strength and the limits of the construction. For $A=k\langle x,y\rangle/(xy,yx)=k[x,y]/(xy)$,
\[
 C_A(w_x,w_y)=\frac{1-w_xw_y}{(1-w_x)(1-w_y)},\quad
 C_A(z)=\frac{1+z}{1-z},\quad E_A(z)=\frac{1-z}{1+z}.
\]
The completed order $\Lambda=k[[x,y]]/(xy)$ has normalization
$B=k[[x]]\times k[[y]]$ and conductor square
\[
 \Lambda=B\times_{k\times k}k,
\]
where the lower map is diagonal and $B\to k\times k$ evaluates the two origins. Equality of constant terms proves this description: $(f,g)$ corresponds uniquely to $f(x)+g(y)-f(0)$. No residue map $k((x))\to k$ is involved; $x$ is a unit in that field.

Its simple $k$ has resolution
\[
 \cdots\xrightarrow{\mathrm{diag}(x,y)}\Lambda^2
 \xrightarrow{\mathrm{diag}(y,x)}\Lambda^2
 \xrightarrow{[x\ y]}\Lambda\longrightarrow k\longrightarrow0.
\]
Indeed $\operatorname{ann}(x)=(y)$, $\operatorname{ann}(y)=(x)$ and $(x)\cap(y)=0$. With $t=x+y$, $\Lambda$ is free of rank two over $k[[t]]$, whereas $k$ is $t$-torsion. Positive syzygies are torsion-free over this DVR. Thus the simple is not periodic from degree zero. Moreover $(x)\oplus(y)$ and $(y)\oplus(x)$ are isomorphic: the unlabelled syzygy tail has period one even though the ordered branch pattern has period two. Tensoring with $k[[t]]/(t)$ does not preserve this resolution: the special fiber is $k[x]/(x^2)$, whose simple has Betti numbers $1,1,1,\ldots$, rather than $1,2,2,\ldots$. The missing homology is a Tor contribution.

For the following closed split subclass, require exactly two incoming and two outgoing arrows at each vertex, unique permitted and forbidden successors and predecessors, and two matched permitted-cycle incidences at each original vertex, with no unmatched endpoints. One can then specify hereditary cycle orders $H_c$ over $k[[x_c]]$ and impose equality of the paired diagonal residues. Map each permitted cycle to successive matrix units, with its wrap arrow multiplied by $x_c$. Distinct permitted paths have independent coefficients in these entries, and every tuple with matching diagonal residues is reconstructed coefficient by coefficient. This proves the completed path-algebra/pullback comparison in this split setting. A common central parameter is $t\mapsto(x_c1)_c$; long permitted paths show that the arrow-adic and $t$-adic filtrations are cofinal. Broader Dedekind bases need their own residue identifications, local uniformizers and projective-cover hypotheses. Hereditary normalization does not make the glued order hereditary.

The derived category $D(A)$ is obtained from complexes of right $A$-modules by inverting quasi-isomorphisms. Its perfect subcategory $\operatorname{Perf}(A)$ consists of objects represented by bounded complexes of finitely generated projectives. A simple with an infinite minimal resolution need not be perfect. The homotopy category of bounded projectives embeds as $\operatorname{Perf}(A)$; $D^b(\mathrm{mod}\text{-}A)$ can be larger. This is the correct finiteness test before forming finite determinants.

For finite-dimensional graded gentle algebras, the surface models of \cite{OPS,Chang,LZZ,Koszul} describe string objects by graded curves and band objects by closed curves with local systems. A line field/degree convention and endpoint or puncture behavior are essential. Appropriate oriented graded intersections give morphism bases, subject to exceptional closed-curve identity and Auslander--Reiten maps. Thus a bare count of transverse intersections is not a universal Hom formula. In a module heart, $\operatorname{Ext}^n(M,N)=\operatorname{Hom}_{D^b}(M,N[n])$; resolutions compute these groups and derived tensor products compute Tor. Rotating or zigzagging curves supplies concrete projective covers in the stated models. Koszul duality exchanges permitted and forbidden combinatorics through a derived functor with finiteness/completion conditions; a formal matrix inverse alone is not a derived equivalence. There is no justified transfer from these finite-dimensional categories to every completed arithmetic order, nor an identification of curve length with prime norm.

\section{Determinant lines, chain endomorphisms and zeta functions}
For a bounded complex $P^\bullet$ of finite projective modules over a commutative base $R$, its determinant is the graded invertible module
\[
 \det(P^\bullet)=\bigotimes_i(\det_R P^i)^{(-1)^i}.
\]
It is a line, with the coherent signs of the determinant functor, not a scalar analytic function. Quasi-isomorphism invariance and derived base change hold in the perfect category \cite{KM,Breuning,DetRecent}. For noncommutative $A$ first choose a realization to a perfect complex over a commutative base; arbitrary $A$-projectives have no ordinary top-exterior-power determinant of this kind. A resolution is a computational presentation, not an additional arithmetic eigenvalue source.

\begin{theorem}[Cohomological cancellation]\label{thm:cohom}
Let $C^\bullet$ be a bounded finite-dimensional complex over a field and let $F$ be a chain endomorphism. Then
\[
 Z(C,F;u):=\prod_i\det(1-uF\mid C^i)^{(-1)^{i+1}}
  =\prod_i\det(1-uF\mid H^i(C))^{(-1)^{i+1}}.
\]
An acyclic complex contributes one.
\end{theorem}
\begin{proof}
Set $B^i=\operatorname{im}d^{i-1}$ and $Z^i=\ker d^i$. Commutation with $d$ preserves both. The two exact sequences
$0\to Z^i\to C^i\to B^{i+1}\to0$ and
$0\to B^i\to Z^i\to H^i\to0$ make $F$ block triangular in adapted bases. Hence its characteristic determinant on $C^i$ is the product of those on $B^i,H^i,B^{i+1}$. The two boundary factors have opposite adjacent exponents and cancel. The finite product requires no convergence argument.
\end{proof}

An infinite periodic resolution is not a bounded complex for this theorem. A truncation leaves a boundary syzygy; omitting it changes the object. Likewise the matrices $U$ and $-V$ in Theorem~\ref{thm:cartan} can package a formal virtual determinant, but have not been shown to be an arithmetic endomorphism commuting with the resolution differential. These are separate constructions.

Over a characteristic-zero field (or after extending scalars to a $\mathbb Q$-algebra), for a finite matrix $M$, differentiation of the determinant at the identity and the formal inverse series give
\[
 -\log\det(1-uM)=\sum_{r\geq1}\frac{\tr(M^r)}r u^r.
\]
Indeed the derivative of the left side is $\tr(M(1-uM)^{-1})$; expand and integrate term by term with constant term zero. For a partial permutation this becomes the primitive-cycle Euler product. Such transfer zeta functions are not automatically Hasse--Weil or number-field $L$-functions. An infinite-dimensional extension needs a stated Fredholm or regularized determinant and its convergence theorem.

\subsection{Testing a positive Hodge replacement}
A natural next attempt is to turn the resolution differential into a positive operator. The following calculation identifies exactly what this retains.
\begin{theorem}[Finite Hodge determinant cancellation]\label{thm:hodge-cancel}
Let $C^\bullet$ be a bounded finite-dimensional complex Hilbert complex, $D=d+d^*$ and $\Delta=D^2$. Write $E,O$ for the even and odd sums, of dimensions $e,o$. Then
\[
 \det(1+z\Delta_E)=\det(1+z\Delta_O),\qquad
 \mu^o\det(\mu1_E+\Delta_E)=\mu^e\det(\mu1_O+\Delta_O).
\]
The shifted parity quotient is $\mu^{e-o}$ as a meromorphic function; pointwise division requires a nonzero denominator. Thus every nonzero spectral factor cancels in the plain alternating Hodge determinant.
\end{theorem}
\begin{proof}
In orthonormal bases the odd self-adjoint $D$ has block form
$\left(\begin{smallmatrix}0&B^*\\B&0\end{smallmatrix}\right)$, so $\Delta_E=B^*B$ and $\Delta_O=BB^*$. Sylvester's identity applied to $zB^*,B$ proves the first assertion. For $\mu\ne0$, factor $\mu$ from each matrix and apply that identity at $z=\mu^{-1}$. The cleared equality is polynomial in $\mu$, hence also holds at zero. The difference $e-o$ is the Euler characteristic, by cancellation of boundary dimensions. The quotient statement is meromorphic; it does not assign a pointwise value to a common $0/0$ at a negative Laplace eigenvalue.
\end{proof}
The ordinary determinant avoids this cancellation but is metric dependent even on the fixed exact complex $\mathbb C\xrightarrow{1}\mathbb C$: metrics $1,r>0$ give both Laplacians $r$, with ordinary factor $r-z$. Also $d(s)=s$ gives $\Delta(s)=|s|^2$, which is not holomorphic near $s=1$. Positivity of a Hodge construction therefore does not supply a determinant-preserving arithmetic comparison. Degree-weighted analytic torsion \cite{RaySinger} is a different invariant and is not excluded by this unweighted parity calculation.

\subsection{A relative boundary realization and its limit obstruction}
\begin{proposition}[Relative realization of the Cartan determinant]\label{prop:boundary}
Let $S$ be the unilateral shift on $\ell^2(\mathbb N)$ and $P_0$ its initial rank-one projection. For a finite matrix $M$ with $1-M$ invertible, define
\[
 A_M=1-(S^*S)\otimes M,\qquad C_M=1-(SS^*)\otimes M.
\]
Then $A_MC_M^{-1}=1-P_0\otimes M$, a finite-rank perturbation of the identity, and
\[
 R(M):=\det_F(A_MC_M^{-1})=\det(1-M).
\]
Consequently $R(-V)/R(U)=\det C_A$ for the arrow-state matrices above, initially where the operators are defined and then by the indicated polynomial/meromorphic continuation.
\end{proposition}
\begin{proof}
Coordinates give $S^*S=1$, $SS^*=1-P_0$. On $P_0\otimes\mathbb C^r$, $C_M$ is the identity and $A_M=1-M$; on its orthogonal complement both equal $1\otimes(1-M)$. Their relative product therefore has the stated form. Its Fredholm determinant is the finite determinant on its $r$-dimensional support. Theorem~\ref{thm:cartan} gives the last assertion. Individual Fredholm determinants of the bulk operators are generally undefined; no quotient of two such determinants is taken. At singular $1-M$ only the scalar polynomial continuation is asserted, not a defined quotient operator.
\end{proof}
This is an explicit Hilbert-space relative realization of the formal Cartan invariant, not an arithmetic or generally self-adjoint realization. Its natural finite sections expose a further analytic gap. For scalar $M=1/2$, the truncated shift on $N\ge2$ coordinates gives a relative operator equal to $1/2$ at the first coordinate, $2$ at the last and $1$ elsewhere. Its determinant is always $1$. Extending it by the identity, it converges strongly to $1-P_0/2$, whose Fredholm determinant is $1/2$. The difference from that limit is the last-coordinate projection, of trace norm exactly $1$. Hence strong convergence here cannot justify determinant convergence. The extended edition gives the general matrix formula and the necessary infinite-dimensional qualifications.

\section{What finite-field purity supplies}
For a separated finite-type $X/\mathbb F_q$ and a suitable constructible $\ell$-adic coefficient system $\mathcal F$, $\ell\ne\operatorname{char}\mathbb F_q$, the Grothendieck trace formula identifies
\[
 L(X,\mathcal F,u)=\prod_{x\ {\rm closed}}
 \det(1-u^{\deg x}\operatorname{Fr}_x\mid\mathcal F_x)^{-1}
 =\prod_i\det(1-u\operatorname{Fr}\mid H_c^i(X_{\overline{\mathbb F}_q},\mathcal F))^{(-1)^{i+1}}.
\]
We choose the Frobenius convention giving eigenvalue $q$ on $H^2(\mathbb P^1)$. For a smooth projective geometrically connected curve and constant coefficients this is
\[
 Z(X,u)=\frac{\det(1-u\operatorname{Fr}\mid H^1)}{(1-u)(1-qu)}.
\]
Deligne's purity theorem gives every complex absolute value $\sqrt q$ on $H^1$ \cite{WeilI,WeilII}; with $u=q^{-s}$ the numerator zeros have $\Rea s=1/2$. This paragraph invokes deep external arithmetic theorems, not consequences of the elementary path-resolution proof. Weil II also treats morphisms over $\mathbb Z$ with fiberwise weights; its scope is not confined to one finite field. What is absent here is a global completed number-field determinant and a positivity theorem for its zero-bearing part. Open or singular spaces require mixed weights or intersection complexes: for example $Z(\mathbb G_m,u)=(1-u)/(1-qu)$ has a weight-zero zero.

One useful finite-field test is $X:y^2=x^3-x$ over $\mathbb F_5$. Its eight points give trace $-2$ and
\[
 Z(X,u)=\frac{1+2u+5u^2}{(1-u)(1-5u)},\qquad
 \alpha_1,\alpha_2=-1\pm2i,
\]
with $|\alpha_i|=\sqrt5$. This benchmarks a genuine cohomological determinant. A nodal or gentle category is not automatically smooth and proper over $k$, the hypotheses of modern noncommutative Weil frameworks \cite{Tabuada}. Surface topology alone supplies neither point counts nor this Frobenius realization.

\begin{proposition}[Duality versus positive purity]\label{prop:purity}
An alternating duality $F^tJF=qJ$ does not force eigenvalues of $F$ to have modulus $\sqrt q$. A positive Hermitian form $H$ with $F^*HF=qH$, $q>0$, does force that conclusion.
\end{proposition}
\begin{proof}
For the first assertion take $q=2$, $F=\mathrm{diag}(2,1)$ and $J=\left(\begin{smallmatrix}0&1\\-1&0\end{smallmatrix}\right)$. Direct multiplication gives the duality, but its eigenvalues have squares of absolute values $4$ and $1$. For the second, if $Fv=\alpha v$ and $v\ne0$, then
$|\alpha|^2v^*Hv=(Fv)^*H(Fv)=qv^*Hv$; divide by $v^*Hv>0$.
\end{proof}
Functional equations similarly pair zeros without placing them on a line. Local weight zero for finite-image Artin Frobenius, already known since its eigenvalues are roots of unity, is insufficient for global number-field GRH.

\section{A precise arithmetic program and a closure theorem}
\begin{proposition}[Positive-length Cartan obstruction]\label{prop:nogo}
In Theorem~\ref{thm:cartan}, specialize $w_a(s)=\eta_a e^{-s\lambda_a}$ with $|\eta_a|\leq1$ and $\lambda_a>0$. Its rational Cartan determinant and its reciprocal are holomorphic and nonzero on $\Rea s>0$. Neither, multiplied by holomorphic nonvanishing factors, can converge locally uniformly to a nonzero holomorphic target having a zero there.
\end{proposition}
\begin{proof}
Every primitive cycle has positive total length and product weight of modulus strictly less than one in this half-plane. Each numerator and denominator factor is therefore nonzero; their finite product is holomorphic with holomorphic reciprocal. If a sequence of these functions converged locally uniformly to a nonzero target with a zero, choose a small closed disk around an isolated target zero and no boundary zero. Uniform boundary convergence and Rouch\'e's theorem force an approximant zero in the disk, a contradiction. The identical argument handles the stated normalizations.
\end{proof}
Thus the resolution/Cartan formula rigorously diagnoses a failed direct Euler-product repair. Artificially centering the weights at $s-1/2$ puts their elementary cycle zeros on a chosen line but proves no matching with an arithmetic $L$-function.

\begin{theorem}[Conditional positive determinant closure]\label{thm:closure}
Let $\Omega$ be a connected open subset of $\mathbb C$ containing an open portion of the critical strip, and let $L_\circ$ be holomorphic on $\Omega$, not identically zero. Suppose holomorphic functions $D_N$ converge locally uniformly to $L_\circ$ and every zero of $D_N$ in $\Omega$ lies on $\Rea s=1/2$. Then every zero of $L_\circ$ in $\Omega$ lies on that line. One sufficient finite realization is
\[
 D_N(s)=g_N(s)\det\bigl(A_N-[-i(s-1/2)]B_N\bigr),
 \quad A_N=A_N^*,\quad B_N>0,
\]
with $g_N$ holomorphic and nonzero and the same convergence hypothesis.
\end{theorem}
\begin{proof}
At a singular pencil, a nonzero $v$ satisfies $A_Nv=zB_Nv$. The quotient $z=(v^*A_Nv)/(v^*B_Nv)$ is real since its denominator is positive. For $z=-i(s-1/2)$ this means $\Rea s=1/2$. If the limit had an off-line zero, a small disk avoiding the line and other zeros would contradict Rouch\'e exactly as in Proposition~\ref{prop:nogo}.
\end{proof}
This theorem is a closure mechanism, not a construction of $A_N,B_N,g_N$. Our existing arithmetic partial-sum approximants converge with compact error bounds; no positive-pencil determinant identity for them has been proved. Replacing their matrices by Hermitian parts can change the determinant and is not such an identity.

A complete attempt using resolutions must discharge the following sequence of concrete claims.
\begin{enumerate}
\item Construct the arithmetic category and realization independently of unknown zeros. Specify coefficients, ramification, actual Galois/Frobenius action, archimedean data and conductor. A dessin's topological monodromy is not this action on one fixed algebra.
\item Choose a finite perfect realization, or an infinite analytic complex with a justified determinant. Prove chain-map compatibility and comparison through conductor gluing and derived base change; include boundary and singularity contributions. Retain determinant trivializations and normalization.
\item Prove the completed $L$-function identity, including gamma factors and multiplicities. Finite-dimensional finite-field rationality in $u$ does not provide a finite-dimensional global Frobenius for primes of different norms. Number-field flow proposals \cite{Deninger} require a regularized determinant and a global generator.
\item Construct and prove positivity/polarization for the \emph{same} zero-bearing realization, or prove a determinant-preserving positive-pencil comparison for the convergent arithmetic approximants. Alternating duality, local unitarity, positivity of unrelated moments and a formal Euler identity do not discharge this step.
\item Prove trace-norm or compact-uniform determinant convergence on an exhaustion, with a nonzero anchor and appropriate pole removal. Then apply Theorem~\ref{thm:closure}, with each intended $L$-function scope and all exceptional points checked.
\end{enumerate}
The local homological part can now be implemented explicitly. The global arithmetic identity, its positive realization and its analytic comparison remain mathematical research problems. No assumption of GRH, Artin holomorphy or unknown zero locations is used in the proved algebraic statements.

\section{Verification and reproducibility}
The separate pinned Lean4.24/mathlib layer verifies a general cleared version of the rectangular determinant identity, acyclic-pair characteristic-determinant cancellation, the actual alternating-duality counterexample, and exactness/projectivity components of a dual-number resolution over an actual nonsemisimple algebra. The Hodge module also verifies actual block squaring, Gram positivity and determinant cancellation, positive scalar metrics and their moving-zero example, and a universal ring identity for the boundary relative product. The infinite Hilbert-space/Fredholm and trace-norm arguments just given remain prose. Existing files verify positive Hermitian pencils and analytic nonvanishing components. Theorems print their types and foundation dependencies. These are components, not a formalization of all forbidden-chain resolutions, the derived surface models, the trace formula, Deligne purity, or GRH.

A standard-library exact rational checker tests finite quivers, path/chain coefficients, differentials and bounded internal-degree exactness, signed matrix inversion and the weighted cycle determinant, including nodal, terminating, periodic and cancellation examples. Finite tests supplement the universal proofs; they do not certify the arithmetic bridge. All source submodules and pinned environments are preserved. Both editions compile from shared core text; the extended edition adds explicit computations, geometric hypotheses and implementation tasks.

\begin{thebibliography}{99}
\bibitem{S1} A. Schreiber, \emph{Surface Algebras I: Dessins D'enfants, Surface Algebras, and Dessin Orders}, \href{https://arxiv.org/abs/1810.06750v5}{arXiv:1810.06750v5} (2023).
\bibitem{S2} A. Schreiber, \emph{Surface Algebras and Surface Orders II: Affine Bundles on Curves}, \href{https://arxiv.org/abs/1812.00621v1}{arXiv:1812.00621v1}; record2018, downloaded PDF internally dated November27,2024.
\bibitem{S3} A. Schreiber, \emph{Multiparameter Persistent Homology---Generic Structures and Quantum Computing}, \href{https://arxiv.org/abs/2210.11433v1}{arXiv:2210.11433v1} (2022); detailed full-reading review in the project archive.
\bibitem{BH} C. Bessenrodt and T. Holm, \emph{Weighted locally gentle quivers and Cartan matrices}, \href{https://arxiv.org/abs/math/0511610v1}{arXiv:math/0511610v1} (2005), Theorem~1 and Proposition~3.4.
\bibitem{FOZ} S. Ford, A. Oswald and J.J. Zhang, \emph{Homological conditions on locally gentle algebras}, \href{https://arxiv.org/abs/2409.08333v3}{arXiv:2409.08333v3} (September2026), \S4.1, Lemma~4.5.
\bibitem{OPS} S. Opper, P.-G. Plamondon and S. Schroll, \emph{A geometric model for the derived category of gentle algebras}, \href{https://arxiv.org/abs/1801.09659v7}{arXiv:1801.09659v7} (July2025), Theorem~3.3 and its exceptions.
\bibitem{Chang} W. Chang, \emph{Geometric models for the algebraic hearts in the derived category of a gentle algebra}, \href{https://arxiv.org/abs/2305.03998v2}{arXiv:2305.03998v2} (2024), Theorem~2.16, Proposition~2.23 and Corollary~2.28.
\bibitem{LZZ} Y.-Z. Liu, C. Zhang and H. Zhang, \emph{Constructing projective resolution and taking cohomology for gentle algebras in the geometric model}, \href{https://arxiv.org/abs/2308.07220v1}{arXiv:2308.07220v1} (2023).
\bibitem{Koszul} Z. Li, Y. Qiu and Y. Zhou, \emph{A geometric realization of Koszul duality for graded gentle algebras}, \href{https://arxiv.org/abs/2403.15190v1}{arXiv:2403.15190v1} (2024), body Theorems~3.6 and3.9.
\bibitem{Cuts} W. Chang, H. Jin and S. Schroll, \emph{Recollements of partially wrapped Fukaya categories and surface cuts}, \href{https://arxiv.org/abs/2206.11196v2}{arXiv:2206.11196v2} (2023).
\bibitem{ResolveTrees} B. Dequ\^ene and M. Schoonheere, \emph{Resolving subcategories for gentle algebras I: Monogeneous resolving subcategories for gentle trees}, \href{https://arxiv.org/abs/2502.20994v3}{arXiv:2502.20994v3} (2026); and PartII, \href{https://arxiv.org/abs/2509.10120v2}{arXiv:2509.10120v2} (2025).
\bibitem{BE} D.A. Buchsbaum and D. Eisenbud, \emph{What makes a complex exact?}, J. Algebra25 (1973),259--268, criterion pp.260--261, \href{https://eisenbud.github.io/papers/pdfs/1973-002.pdf}{primary paper}.
\bibitem{Bruns} W. Bruns, \emph{The existence of generic free resolutions and related objects}, Math. Scand.55 (1984),33--46, Theorem~1.
\bibitem{Eisenbud} D. Eisenbud, \emph{Homological algebra on a complete intersection, with an application to group representations}, Trans. AMS260 (1980),35--64, \S\S5--6.
\bibitem{BL} A. Beauville and Y. Laszlo, \emph{Un lemme de descente}, C. R. Acad. Sci. Paris320 (1995),335--340, \href{https://math.univ-cotedazur.fr/~beauvill/pubs/descente.pdf}{author text}.
\bibitem{Magyar} P. Magyar, \emph{Affine Schubert varieties and circular complexes}, \href{https://arxiv.org/abs/math/0210151v1}{arXiv:math/0210151v1} (2002).
\bibitem{KM} F.F. Knudsen and D. Mumford, \emph{The projectivity of the moduli space of stable curves I}, Math. Scand.39 (1976),19--55, \href{https://doi.org/10.7146/math.scand.a-11642}{doi:10.7146/math.scand.a-11642}.
\bibitem{Breuning} M. Breuning, \emph{Determinant functors on triangulated categories}, \href{https://arxiv.org/abs/math/0610435}{arXiv:math/0610435} (2006).
\bibitem{DetRecent} D. Eriksson and G. Freixas i Montplet, \emph{Deligne--Riemann--Roch and intersection bundles}, J. \'Ec. polytech. Math.11 (2024),247--361, \S4.5, \href{https://doi.org/10.5802/jep.254}{doi:10.5802/jep.254}.
\bibitem{WeilI} P. Deligne, \emph{La conjecture de Weil I}, Publ. Math. IH\'ES43 (1974),273--307, \S1, especially Theorem~1.6 and equations1.5.1--1.5.4, \href{https://www.numdam.org/item/PMIHES_1974__43__273_0.pdf}{primary text}.
\bibitem{WeilII} P. Deligne, \emph{La conjecture de Weil II}, Publ. Math. IH\'ES52 (1980),137--252, Theorem~3.3.1 and Corollaries~3.3.4--3.3.6, \href{https://numdam.org/articles/10.1007/BF02684780/}{primary text}. See also SGA4\ensuremath{\frac12}, Rapport, \S3, for the trace formula.
\bibitem{Tabuada} G. Tabuada, \emph{Noncommutative Weil conjecture}, \href{https://arxiv.org/abs/1808.00950v7}{arXiv:1808.00950v7}; smooth proper dg categories over finite fields.
\bibitem{Deninger} C. Deninger, \emph{Arithmetic geometry and analysis on foliated spaces}, \href{https://arxiv.org/abs/math/0505354}{arXiv:math/0505354} (2005); conjectural number-field interfaces distinguished from proved function-field results.
\bibitem{RaySinger} D.B. Ray and I.M. Singer, \emph{R-Torsion and the Laplacian on Riemannian manifolds}, Advances in Math.7 (1971),145--210, \S1, Definition1.6 and Proposition1.7, \href{https://doi.org/10.1016/0001-8708(71)90045-4}{doi:10.1016/0001-8708(71)90045-4}.
\bibitem{Bornemann} F. Bornemann, \emph{On the numerical evaluation of Fredholm determinants}, \href{https://arxiv.org/abs/0804.2543v1}{arXiv:0804.2543v1}, \S\S3--5; Math. Comp.79 (2010),871--915. Only its general determinant/trace-class convergence statements are used here.
\end{thebibliography}

\end{document}
